\uuid{yMOe}
\titre{Distance}
\auteur{bodin ; hanani ; mzari}
\organisation{exo7}
\licence{CC BY-NC-SA 4.0}
\datecreate{2020-01-20}
\niveau{L1}
\module{Géométrie}
\chapitre{Géométrie de l'espace}
\sousChapitre{Distance}
\difficulte{1}
\theme{Distance}
\exo7qcmid{lille-1.geometrie-espace.022}

\contenu{
\texte{
Pour ces questions, l'espace est muni d'un repère orthonormé direct  $(O, \vec{i}, \vec{j}, \vec{k})$.
}

\begin{question}
\qclassification{141.05}
\qtags{motcle=distance}
Soit $A(-1,1,0)$ et $P$ le plan passant par $B(1,0,1)$ et dirigé par les vecteurs $\vec{u}(1,1,1)$ et $\vec{v}(1,0,-1)$. La distance de $A$ à $P$ est :
\begin{answers}

    \bad {$\frac{1}{\sqrt 6}$}

    \good {$\frac{5}{\sqrt 6}$}

     \bad{$\sqrt 6$}

    \bad{$\frac{4}{\sqrt 6}$}

\end{answers}
\begin{explanations}
Si $P$ passe par un point $B$ et est dirigé par des vecteurs $\vec{u}$ et $\vec{v}$ et $A$ un point, la distance de $A$ à $P$ est donnée par : $\frac{ \vert \det (\overrightarrow{BA}, \vec{u},\vec{v})\vert}{\Vert \vec{u} \wedge \vec{v}\Vert}$.
\end{explanations}
\end{question}
}
